For a coefficient characteristic $\ell\ne p$, Carayol's published
Theorem~A gives the full local Weil--Deligne compatibility, not only
an equality at unramified primes
\cite[\S\S0.3--0.7, pp.~409--411]{carayol}.
Its base field here is $\mathbb Q$, its real weight is two, and its
extra finite discrete-series-place condition applies only to even
degree totally real fields, so imposes no additional hypothesis
here. Theorem~A includes the dyadic cases; the weaker Theorem~B
is not the result being used.

Carayol normalizes reciprocity by sending geometric Frobenius to
a uniformizer. His Hecke correspondence in \S0.5 differs from
local Langlands by a contragredient and an unramified norm twist.
It is characterized there by equality of local $L$- and
epsilon-factors, including character twists. The conductor is
unchanged by either duality or an unramified twist. In particular,
the passage to the Tate-module convention used for $r_\ell$ changes
no conductor exponent. Equality of good Frobenius traces identifies
the global representation supplied by Carayol with the already
constructed irreducible modular representation, in this convention.
This uses Chebotarev and Brauer--Nesbitt; it does not use the
unpublished irreducibility observation in Carayol's remark after
Theorem~A. Since the matching representation is irreducible, its
matching semisimplification also forces the original representation
to be irreducible.

The epsilon-factor conductor in (CL1) therefore equals the Artin
conductor of $r_\ell|_{G_{\mathbb Q_p}}$:
\[
 a_p(r_\ell)=a(\Pi_{f,p})=v_p(N),\qquad p\ne\ell.
\tag{CL2}
\]
The Artin conductor here includes its tame term and Swan term;
equivalently one uses the full Weil--Deligne conductor, including
the monodromy contribution. One must not discard monodromy by
replacing the local representation with its semisimplification on
inertia. Carayol's compatibility supplies precisely the required
full local object. His epsilon-factor normalization and
Casselman's conductor-zero additive character give the same
exponent; changing an additive character only introduces the usual
common dimension-dependent shift, so no such change is needed.

Apply (CL2) with $\ell=5$ at every $p\ne5$ and with $\ell=3$ at
$p=5$. The companion proof identifies both $r_5$ and $r_3$ with
the realizations of this same $f$. Thus the asserted conductor
bounds control every prime divisor of the primitive level, even
when $p^2\mid N$. No compatibility theorem at $p=\ell$ is needed.
